\scanpage{55}
\mspage{54}

\uline{Scholie} Se donner un atelier (quasi) ensembliste, ``revient au $=$'' que de se donner un couple $(L,\unsure{\mathcal F|\mathscr F})$, où $L$ est un ens, $\mathcal F$ un sous-atelier de $\operatorname{Fig}(L)$, satisfaisant \unsure{la|les} condition\deleted{\unsure{s}} suivante\deleted{\unsure{s}} :

\deleted{\uline{At ens$^{*}$ 1} \deleted{\unsure{(sé)|(sép.)|(*)}} Les strates des $\unsure{\Phi|\underline{\Phi}}\in\mathcal F$ séparent les points de $L$, i.e. si $x,y\in L$, $x\ne y$, alors $\exists\,\deleted{\unsure{\Phi}}\unsure{P|F}\in\Phi$, $A\in \unsure{P|F}$, tel qu'on ait $x\in A$, $y\notin A$ ou $y\in A$, $x\notin A$.}

\[
\text{\uline{At ens}}^{*}\,\deleted{\unsure{2}}1\quad \deleted{\unsure{\varnothing .}}\quad
\forall\,x\in L,\quad\{\{x\}\}\in\mathcal F
\]

\deleted{\unsure{Le}} Le tout est de vérifier que, dans ce cas, $L$ se réconstitue à partir de $(\mathcal F,\leq,\ll)$ comme l'ens des lieux de $\mathcal F$, donc que
\[
x\longmapsto\{\{x\}\}=\Phi_x
\]
est une bijection de $L$ \deleted{\unsure{(ou de $M$}}\inserted{avec} l'ens des éléments de $\mathcal F\setminus\{\varnothing\}$ minimaux pour $\ll$. \inserted{Il est clair que $\Phi_x$ est minimal dans $\operatorname{Fig}(L)\setminus\unsure{\{\varnothing\}|(\varnothing)}$, a fortiori dans $\mathcal F\setminus\{\varnothing\}$}.

Or soit $F\in\deleted{\unsure{\Phi}}\mathcal F\setminus\{\varnothing\}$, \inserted{minimal} \deleted{Soit $A\in F$, on a $A^{\circ}\ne\varnothing$, soit $x\in A^{\circ}$ alors $\Phi$} \deleted{Soit} On a $|F|\ne\varnothing$ (puisque $F\ne\varnothing$) soit $x\in |F|$, alors $\Phi_x\ll F$, donc $\Phi_x=F$, cqfd.
